% Figure for Classical Mechanics §B5.2: three wires from A = (0,0) to B = (pi a, 2a), y measured downward, g = a = 1; descent times computed (cycloid exactly pi).
\begin{tikzpicture}
  \begin{axis}[nfig, width=10.5cm, height=8.2cm, y dir=reverse, xmin=0, xmax=3.8, ymin=0, ymax=3.3,
      axis x line=top, axis y line=left, xlabel={$x/a$}, ylabel={$y/a$ (down)},
      xtick={1,2,3}, ytick={1,2}, clip=false, axis equal image,
      legend style={at={(axis cs:3.78,3.28)}, anchor=south east, font=\small, draw=black!30, fill=white, row sep=1pt},
      legend cell align=left]
    % straight line
    \addplot[domain=0:1, samples=2, figR, line width=1.1pt] ({pi*x}, {2*x});
    \addlegendentry{straight line, $T = 3.724\sqrt{a/g}$}
    % circular arc with vertical tangent at A: centre (c,0), c = (pi^2+4)/(2 pi)
    \addplot[domain=0:115.036, samples=80, figG, line width=1.1pt] ({2.2074161 - 2.2074161*cos(x)}, {2.2074161*sin(x)});
    \addlegendentry{circular arc, $T = 3.221\sqrt{a/g}$}
    % cycloid
    \addplot[domain=0:180, samples=90, figB, line width=1.4pt] ({rad(x) - sin(x)}, {1 - cos(x)});
    \addlegendentry{cycloid, $T = \pi\sqrt{a/g} = 3.142\sqrt{a/g}$}
    % rolling circle generating the cycloid (theta = 1.3)
    \addplot[domain=0:360, samples=90, black!40, dashed, line width=0.5pt, forget plot] ({1.3 + cos(x)}, {1 + sin(x)});
    \fill[black!45] (axis cs:1.3,1) circle (1.2pt);
    \draw[black!45, thin] (axis cs:1.3,1) -- (axis cs:{1.3-sin(deg(1.3))},{1-cos(deg(1.3))});
    \fill[figB] (axis cs:{1.3-sin(deg(1.3))},{1-cos(deg(1.3))}) circle (1.6pt);
    \node[black!55, font=\scriptsize, anchor=west] at (axis cs:2.33,1.0) {rolling circle};
    \fill (axis cs:0,0) circle (2pt) node[below left] {$A$};
    \fill (axis cs:pi,2) circle (2pt) node[right, xshift=1pt] {$B$};
  \end{axis}
\end{tikzpicture}
